\section{Introduction}

High-speed atomic force microscopy (HS-AFM) enables direct visualization of biomolecular dynamics in solution under near-physiological conditions, providing unprecedented insights into the relationship between biomolecular conformational dynamics and biological function \cite{ando_high-speed_2008, ando_high-speed_2022}. 
However, AFM measurements are inherently limited to surface topography, and the achievable spatial resolution is fundamentally constrained by the finite size of the cantilever tip \cite{villarrubia_algorithms_1997, matsunaga2023endtoend}. 
To address these limitations, there is a compelling need for computational modeling approaches that can integrate experimental AFM data with atomic-resolution three-dimensional structural models.

One approach is rigid-body fitting, in which template structures are selected from structural databases or molecular dynamics (MD) simulations and then exhaustively rotated and translated to find the best pose matching the AFM image \cite{ando_high-speed_2005, scheuring_structure_2005, scheuring_high-resolution_2007,asakawa_submolecular-scale_2011, trinh_computational_2012, chaves_conformational_2013, amyot_bioafmviewer_2020, amyot_simulation_2022, niina_rigid-body_2021}. 
Although powerful, this approach cannot account for conformational flexibility or changes; consequently, rigid-body fitting often fails to provide accurate structural models when the template library lacks conformations that match the AFM image. 
In real biological systems, the function of biomolecules is often coupled with conformational flexibility or dynamics. 
Thus, methods accounting for such conformational flexibility are required for analyzing and interpreting experimental AFM data.

In recent years, methods have been proposed that utilize MD or Monte Carlo (MC) simulation to generate conformations consistent with AFM images. 
For example, in flexible fitting \cite{Niina2020, ngo2025atomic}, the agreement with an AFM image is defined as an additional potential energy function $V_{\mathrm{AFM}}$, and MD simulations are driven to minimize it. 
In flexible fitting, sampling is often trapped in local minima, making optimization slow when separated by high barriers. 
If the global minimum is desired, prohibitively long simulations must be performed to ensure adequate sampling of the conformational space. 
Normal-mode-based approaches achieve substantially faster per-frame processing on standard hardware. 
NMFF-AFM \cite{Wu2024, amyot_flexible_2025} employs an iterative scheme in which the linear normal-mode response is consecutively applied to gradually changing conformations, enabling the reconstruction of large-amplitude anharmonic transitions. 
AFMFit \cite{Vuillemot2025AFMfit} takes a different strategy, combining FFT-based pose search with a nonlinear normal-mode method to avoid structural distortions inherent in linear extrapolation and to efficiently process hundreds of molecular images into a conformational ensemble.
However, these approaches rely on MD force fields or elastic-network-derived normal-mode bases, so the accessible conformational space is shaped by the choice of model parameterization and the number of modes retained.
Another approach is to use a 3D Gaussian mixture model combined with MC sampling to estimate 3D low-resolution protein structures with conformational flexibility \cite{dasgupta_reconstruction_2020, dasgupta_reconstruction_2021}.

In the field of AI-based structure modeling or prediction, sequence-to-structure models such as \AFii \cite{jumper2021highly, ahdritz2024openfold, mirdita2022colabfold}, \AFiii \cite{Abramson2024}, and \Boltzi \cite{wohlwend2024boltz1} have demonstrated unprecedented precision in predicting the 3D fold of a protein. 
These models can be regarded as generative AI models that generate atomistic conformational ensembles from a given sequence. 
In principle, these generative AI models can generate not only specific peak structures present in the training data, but also diverse structural conformations. 
Moreover, it is known that the generation process can be navigated by applying conditioning or external forces during the generation stage. 
Recently, guiding the generative process of diffusion-based structure prediction models, including \AFiii and \Boltzi, has emerged as an approach aimed at generating targeted conformational ensembles. 
For example, in the context of NMR and X-ray crystallography data analysis \cite{maddipatla2025inverse}, it has been demonstrated that imposing guided restraints during the generative process enables the generation of ensembles consistent with experimental observations. 
Similarly, for cryo-EM data \cite{raghu2025multiscale}, the incorporation of global or local density-map fitting as restraints has been proposed as a means of navigating the predicted conformations toward agreement with the experimental data.

In this study, we introduce \Model, which extends these ideas to AFM data analysis by employing structural restraints extracted from AFM images to guide the sampling trajectories of pre-trained protein structure generative models. Specifically, in the first step, a group-equivariant convolutional neural network (CNN) \cite{e2cnn} estimates low-dimensional CVs (such as inter-domain distances) from AFM images, and then these estimated values are used as restraints to guide the generative process, enabling the rapid sampling of 3D conformations consistent with AFM images. We validated \Model in twin experiments. First, using adenylate kinase (\AK, PDB IDs: 1AKE \cite{muller1992adenylate} and 4AKE \cite{muller1996adenylate}), we estimated the underlying CV values and accurately reconstructed conformations close to the ground truth. 
Second, we applied \Model to real experimental HS-AFM data from the flagellar protein \flhac (PDBID: 3A5I \cite{saijo-hamano2010flha}) and 
confirmed that it predicts conformations more consistent with AFM images than the crystal structure fitted by rigid-body fitting calculations. 
Furthermore, we confirmed that \Model captures the time-correlated domain motions that reflect underlying conformational dynamics.

A distinctive feature of our framework, \Model, is that all steps—from training data generation, model training, and inference, to structure generation—are integrated within a single unified pipeline, without relying on time-consuming processes such as molecular dynamics simulations. 
This enables rapid reconstruction of structures from AFM images while reducing computational cost, allowing the observation of structural dynamics from hundreds of consecutive HS-AFM images. \Model provides a framework for observing conformational transitions and functional mechanisms underlying HS-AFM images.

\begin{figure}[htbp]
    \centering
    \includegraphics[width=0.75\linewidth]{images/figure-1.pdf}
    \caption{
        \textbf{Schematic overview of \Model and outcome}.
        \textbf{(a)} The \Model framework consists of two phases.
        \textit{Training phase} (top): target CV values are placed on a grid in the CV space and used to navigate \AFiii, generating a diverse set of conformations. Pseudo-AFM images are then rendered from each conformation under uniformly sampled 3D orientations, and the resulting (image, CV) pairs serve as the training data for the g-CNN.
        \textit{Prediction phase} (bottom): the pretrained g-CNN predicts CV values from an input AFM image; during the generative process of \AFiii, where the diffusion step $i$ progresses from 0 to 200, the MSE between the predicted CVs and those of the generated conformation is computed when $i > 108$, and its gradient is added to the \AFiii score to guide the generation toward a conformation consistent with the AFM image.
        \textbf{(b,\,c)} Results of structure prediction from pseudo-AFM images of \AK.
        In (b), the upper panels show the reference pseudo-AFM images generated from the crystal structures (closed: 1AKE, open: 4AKE). The lower panels show the pseudo-AFM images generated from the structures estimated by \Model.
        In (c), the blue structures represent the reference conformations, and the red structures are the corresponding \Model estimates.
        }
    \label{fig:overview}
\end{figure}
